\documentclass[11pt,reqno]{amsart}

\usepackage{amsmath}
\usepackage{amssymb}
\usepackage{amsthm}
\usepackage[dvipsnames]{xcolor}
\usepackage{hyperref}
\usepackage{booktabs}
\usepackage{array}
\usepackage{enumitem}
\usepackage{microtype}

\hypersetup{
  colorlinks=true,
  linkcolor=MidnightBlue,
  citecolor=MidnightBlue,
  urlcolor=MidnightBlue
}

\newtheorem{theorem}{Theorem}[section]
\newtheorem{proposition}[theorem]{Proposition}
\newtheorem{lemma}[theorem]{Lemma}
\newtheorem{corollary}[theorem]{Corollary}
\newtheorem{conjecture}[theorem]{Conjecture}
\theoremstyle{definition}
\newtheorem{definition}[theorem]{Definition}
\newtheorem{remark}[theorem]{Remark}
\newtheorem{example}[theorem]{Example}

\DeclareMathOperator{\disc}{disc}
\DeclareMathOperator{\Spec}{Spec}
\DeclareMathOperator{\Stokes}{St}

\title[CM Predicate for Degree-2 PCFs]%
{Complex Multiplication as a Transcendence Predicate\\
for Degree-Two Polynomial Continued Fractions}

\author{Papanokechi}
\address{Independent researcher, Yokohama, Japan}
\email{}
\thanks{ORCID:
\href{https://orcid.org/0009-0000-6192-8273}{0009-0000-6192-8273}.
The author received no external funding for this work.
The numerical experiments were carried out using the AEAL/SIARC
multi-agent computational methodology described in
\cite{P8AEAL}; AI assistants (GitHub~Copilot powered by Anthropic
Claude, and Anthropic Claude directly) were used to generate code
for high-precision computations and to cross-check intermediate
results. All conjectures, proofs, evidence-class assignments, and
editorial decisions are the sole responsibility of the human author.
The AI assistants are not authors. Session-level computation logs
and prompts are archived in the SIARC research workspace and are
available from the author on request, in line with the AEAL
reproducibility requirement.}

\subjclass[2020]{Primary 11J70, 11J81, 11Y65;
Secondary 34M55, 11G15, 33E17.}

\keywords{Polynomial continued fractions, complex multiplication,
balanced discriminant, Heun equation, Painlev\'e transcendents,
Stokes phenomenon, integer relation detection, experimental
mathematics.}

\date{May 1, 2026}

\begin{document}

\begin{abstract}
We propose a transcendence predicate for degree-two polynomial
continued fractions (PCFs), based on the sign of the
\emph{balanced discriminant}
$\Delta = \beta^{2} - 4\alpha\gamma$ of the denominator polynomial
$b_{n} = \alpha n^{2} + \beta n + \gamma$. Working inside the
$\Spec(K)$ classification framework, we present a v5 upgrade of
the schema that adds the Stokes exponent and the
connection-coefficient proxy as components~6 and~7. Our central
empirical contribution is a sharp dichotomy across thirty
degree-two families: twenty-four families in the $F(2,4)$
Trans-stratum have $\Delta\in\{+1,+8\}$ and admit elementary
closed forms, while six candidates with
$\Delta\in\{-3,-4,-7,-11,-20,-28\}$ all return PSLQ non-detection
against an eighteen-element constants basis at working precision
$220$ digits and integer bound $10^{10}$. We promote these
observations to Conjecture~A v5 (parts~(i)--(iv)). The
Painlev\'e~III$(D_{6})$ transcendental nature of the
$\Delta=-11$ prototype $V_{\mathrm{quad}}$ is recalled, and a
\emph{Stokes-exponent diagnostic} is introduced and verified
across all six $\Delta<0$ degree-two families, including a
non-Heegner intra-field cross-check (QL06 with $\Delta=-7$
and QL26 with $\Delta=-28$, both in $\mathbb{Q}(\sqrt{-7})$).
The intra-field replication identifies the CM field, not the
class number, as the structural carrier of the predicate.
The proven $\Delta>0$ side of the dichotomy is established as
Theorem~\ref{thm:trans-closed}.
\end{abstract}

\maketitle

%==========================================================
\section{Introduction}\label{sec:intro}
%==========================================================

\subsection{The transcendence predicate}
\label{ssec:predicate}

A \emph{polynomial continued fraction} (PCF) is determined by two
sequences $(a_{n})_{n\ge 1}$ and $(b_{n})_{n\ge 0}$ with
$a_{n},b_{n}\in\mathbb{Q}[n]$. Whenever the formal expression
\begin{equation}\label{eq:PCF}
K(a_{n},b_{n})
\;=\;
b_{0} + \cfrac{a_{1}}{b_{1} + \cfrac{a_{2}}{b_{2} + \cdots}}
\end{equation}
converges, its limit
$L = K(a_{n},b_{n}) \in \mathbb{R}\cup\{\infty\}$ is called the
\emph{value} of the PCF. The natural inverse problem,
familiar since Euler and Brouncker, asks: \emph{given $L$ as a
classical constant, what (if any) PCF representations
$L = K(a_{n},b_{n})$ exist with polynomial coefficients?}
The forward problem, which is the subject of this paper, is the
opposite: \emph{given $(a_{n}, b_{n})$ in a structured family,
predict whether $L$ admits a closed form, and quantify the
obstruction when none exists.}

We work throughout with the \emph{degree-two class}
$d=2$ in the standard Spec$(K)$ convention~\cite{Papanokechi2026Spec},
namely
\begin{equation}\label{eq:d2-class}
a_{n} \;=\; \delta\,n + \varepsilon, \qquad
b_{n} \;=\; \alpha\,n^{2} + \beta\,n + \gamma,
\end{equation}
with $\alpha,\beta,\gamma,\delta,\varepsilon\in\mathbb{Z}$ and
$\alpha,\delta\neq 0$. The associated forward Wallis recurrence
$p_{n} = b_{n}p_{n-1} + a_{n}p_{n-2}$,
$q_{n} = b_{n}q_{n-1} + a_{n}q_{n-2}$,
with seeds $(p_{-1},p_{0},q_{-1},q_{0})=(1,b_{0},0,1)$, computes
the convergents $p_{n}/q_{n}\to L$ at exponential precision rate
$\rho = \log_{10}|\Lambda|$, where
$\Lambda$ is the dominant Perron--Poincar\'e root of the
characteristic equation associated with~\eqref{eq:d2-class}.

The central object of this paper is the \emph{balanced
discriminant} of the denominator polynomial $b(x) = \alpha x^{2} +
\beta x + \gamma$,
\begin{equation}\label{eq:Delta}
\Delta \;=\; \beta^{2} - 4\alpha\gamma \;=\; \disc(b).
\end{equation}
The roots of $b(x)$ lie in the imaginary-quadratic field
$\mathbb{Q}(\sqrt{\Delta})$ when $\Delta<0$, and in $\mathbb{Q}$
or a real-quadratic field when $\Delta\ge 0$. We refer to the
former as the \emph{CM (complex-multiplication) case} and the
latter as the \emph{non-CM case}.

\subsection{Empirical dichotomy and the conjecture}

\begin{conjecture}[Conjecture~A, transcendence dichotomy]
\label{conj:A-intro}
Let $K = K(a_{n},b_{n})$ be a degree-two PCF with rational
coefficients as in~\eqref{eq:d2-class}, and let
$\Delta=\beta^{2}-4\alpha\gamma$. Then:
\begin{enumerate}
\item[(i)] if $\Delta>0$, then $L = K(a_{n},b_{n})$ admits an
elementary closed form, expressible in the field generated over
$\mathbb{Q}$ by $\pi$, $e$, $\log\mathbb{Q}^{\times}$,
$\zeta(3)$, $\sqrt{\mathbb{Q}}$, and $\gamma$;
\item[(ii)] if $\Delta<0$, then $L$ admits no such elementary
closed form;
\item[(iii)] in case~(ii), there exists a one-parameter
CM-respecting deformation $K(t)$ of $K$, with $\Delta(t)<0$ for
all $t$ in a neighbourhood of $0$, under which the local
Stokes exponent
$S(t) = \log|L(t)-L(0)|/\log|t|$ satisfies
$\lim_{t\to 0} S(t) \in (0,1)$;
\item[(iv)] in case~(ii), the singular set of the associated
Heun ODE for $K$ lies in $\mathbb{Q}(\sqrt{\Delta})$.
\end{enumerate}
\end{conjecture}

Part~(i) is established as Theorem~\ref{thm:trans-closed} below
for the explicit $F(2,4)$ Trans-stratum sub-class
($\Delta\in\{+1,+8\}$) that motivates the conjecture. Parts~(ii)
and~(iii) are supported by exhaustive PSLQ non-detection and a
new Stokes-exponent diagnostic across six $\Delta<0$ families;
part~(iv) is a structural fact about the quadratic formula and
is proved in passing. The remaining open content of
Conjecture~\ref{conj:A-intro} is the implication
$\Delta<0 \Rightarrow$ \emph{no} closed form, which is a
genuine transcendence statement and remains the principal
mathematical question.

\subsection{Headline contributions}\label{ssec:headline}

The key contributions of this paper are as follows.

\begin{enumerate}[label={\bfseries(C\arabic*)}]
\item \emph{The balanced discriminant as transcendence
predicate.} We propose, formalise, and substantiate
Conjecture~\ref{conj:A-intro}: a single sign condition on
$\Delta$ controls a sharp transcendence dichotomy across
thirty degree-two PCF families. (Section~\ref{sec:dichotomy}.)

\item \emph{Theorem for the $\Delta>0$ side.} For the
$F(2,4)$~Trans stratum with
$\Delta\in\{+1,+8\}$, we give a closed-form proof in
Theorem~\ref{thm:trans-closed} that the value
$L = K(a_{n},b_{n})$ is in the elementary field of
Conjecture~\ref{conj:A-intro}(i). (Section~\ref{sec:thm}.)

\item \emph{The Stokes-exponent diagnostic.}
Definition~\eqref{eq:S} introduces a numerical exponent
$S(t)$ that detects fractional-power Painlev\'e/Stokes
obstructions even when no specific Painlev\'e ansatz fits
to threshold. We verify across all six $\Delta<0$
degree-two families that $S<1$ under at least one
CM-respecting deformation, completing the empirical
support for Conjecture~\ref{conj:A-intro}(iii).
(Section~\ref{sec:stokes}.)

\item \emph{Intra-field replication ruling out the
Heegner-number hypothesis.} The pair (QL06, QL26) lies in
the same CM field $\mathbb{Q}(\sqrt{-7})$ at $\Delta=-7$
($h=1$, Heegner) and $\Delta=-28$ ($h=2$, non-Heegner).
Both confirm $S<1$ with the same Heun-root structure
$(\cdot \pm i\sqrt{7})/4$. Holding the CM field fixed and
varying the class number does not extinguish the Stokes
signal, identifying the CM field --- not the class number ---
as the structural carrier of the predicate.
(Section~\ref{sec:stokes}, \S\ref{ssec:intrafield}.)

\item \emph{Painlev\'e III$(D_{6})$ prototype.} The
$\Delta=-11$ family $V_{\mathrm{quad}}$ provides the explicit
Painlev\'e prototype with parameters
$(\alpha,\beta,\gamma,\delta) = (1/6, 0, 0, -1/2)$, with
Stokes constant computed to eight digits; we recall the
result and locate it inside the present framework.
(Section~\ref{sec:vquad}.)

\item \emph{Reproducible computational methodology.}
Appendix~\ref{app:numerical} documents the high-precision
Wallis recurrence, the eighteen-element PSLQ basis, the
Heun-root computation, and the deformation/derivative
stencils, at a level sufficient for independent
reproduction.
\end{enumerate}

\subsection{Relation to prior work}\label{ssec:prior}

The Spec$(K)$ classification framework was introduced
in~\cite{Papanokechi2026Spec}, motivated by the
\emph{Ramanujan Machine} programme of
Raayoni--Gottlieb--Manor--Pisha--Harris--Mendlovic--Haviv--Hadad--Kaminer
\cite{Raayoni2021}, which produced large catalogues of
PCF representations of classical constants but offered no
\emph{a priori} predicate for which input families would yield
such representations. The dichotomy proposed here is the simplest
non-trivial such predicate, and it specialises in the $\Delta>0$
case to the closed-form structure already implicit in the
$F(2,4)$~Trans table of~\cite{Papanokechi2026F24}.

The CM/non-CM distinction has a long pedigree in transcendence
theory, going back to Schneider, Lang, Bertrand, Wolfart, and
W\"ustholz, where it controls the transcendence of
$\Gamma$-values, periods of elliptic curves with complex
multiplication, and quasi-modular forms. The novelty here is
that the same dichotomy appears at the comparatively elementary
level of two-coefficient continued fractions over $\mathbb{Z}$,
where the obstruction can be detected by integer-relation
algorithms (Ferguson--Bailey~PSLQ \cite{FergusonBailey1992}) at
modest precision.

The analytic side of Conjecture~\ref{conj:A-intro} connects to
the isomonodromic deformation theory of the Heun--Painlev\'e
correspondence \cite{Okamoto1987}; for the
$\Delta=-11$ prototype, Painlev\'e III$(D_{6})$ machinery
\cite{Papanokechi2026Vquad} provides explicit Stokes data.
The companion preprint
\cite{Papanokechi2026Spec} establishes the Poincar\'e
characteristic lemma in the balanced-discriminant
formulation, on which our $\Lambda$-field assignments rest.

\subsection{Organisation}\label{ssec:org}

Section~\ref{sec:background} reviews the necessary background
on PCFs, $\Spec(K)$, CM fields, the Heun equation, and the
Stokes phenomenon. Section~\ref{sec:thm} states and proves
the proven $\Delta>0$ part as Theorem~\ref{thm:trans-closed}.
Section~\ref{sec:dichotomy} presents the thirty-family
evidence table and parts~(i)--(ii) of Conjecture~A v5.
Section~\ref{sec:stokes} develops the Stokes-exponent
diagnostic and its six-family verification, including the
intra-field replication. Section~\ref{sec:vquad} recalls the
$V_{\mathrm{quad}}$ Painlev\'e III$(D_{6})$ prototype.
Section~\ref{sec:methodology} documents the computational
methodology. Section~\ref{sec:open} records open problems.
Appendix~\ref{app:numerical} contains the numerical methodology
in full detail.

%==========================================================
\section{Background}\label{sec:background}
%==========================================================

\subsection{Polynomial continued fractions and Wallis convergents}
\label{ssec:wallis}

For polynomial sequences $a_{n},b_{n}\in\mathbb{Q}[n]$ the
formal continued fraction~\eqref{eq:PCF} is a particular linear
fractional transformation of the truncations
$L_{N} = b_{0} + a_{1}/(b_{1}+\cdots+a_{N}/b_{N})$. The
\emph{Wallis recurrence}
\begin{equation}\label{eq:wallis}
\binom{p_{n}}{q_{n}}
\;=\;
b_{n}\binom{p_{n-1}}{q_{n-1}}
\;+\;
a_{n}\binom{p_{n-2}}{q_{n-2}},
\qquad
\binom{p_{-1}}{q_{-1}}\;=\;\binom{1}{0},\;
\binom{p_{0}}{q_{0}}\;=\;\binom{b_{0}}{1},
\end{equation}
generates the convergents $L_{N} = p_{N}/q_{N}$. When the
sequence $L_{N}$ converges, its limit $L$ is the value of
$K(a_{n},b_{n})$. The convergence rate is governed by the
Perron--Poincar\'e theorem: the dominant root $\Lambda$ of the
limiting characteristic equation controls
$|q_{n+1}/q_{n}|\to|\Lambda|$, and the rate of convergence
$|L-L_{N}|$ to zero is exponential with rate
$\rho=\log_{10}|\Lambda|$ digits per step.

In the degree-two class~\eqref{eq:d2-class}, the
Perron--Poincar\'e characteristic polynomial is
\begin{equation}\label{eq:char}
\Lambda^{2} - \alpha\Lambda - \delta \;=\; 0,
\end{equation}
and $\Lambda$ is the larger root in absolute value;
the field $\Lambda \in \mathbb{Q}(\sqrt{\alpha^{2}+4\delta})$
records the Perron--Poincar\'e structure.

\subsection{The Spec$(K)$ classification}\label{ssec:speck}

The Spec$(K)$ invariant of~\cite{Papanokechi2026Spec} is the
five-tuple
\begin{equation}\label{eq:speck}
\Spec(K) \;=\; (d, \Lambda, \Delta, \rho, \tau),
\end{equation}
with $d$ the dominant growth degree, $\Lambda$ the normalised
Perron--Poincar\'e root, $\Delta$ the balanced discriminant,
$\rho$ the Wallis tail-decay rate, and $\tau$ the arithmetic
sector. We use the v5 upgrade introduced
in~\cite{Papanokechi2026Spec} which fixes two cross-paper
inconsistencies and adds two analytic components:
\begin{itemize}
\item the \emph{Stokes exponent} $S(t)$ of equation
\eqref{eq:S}, recording local power-law behaviour of
$|L(t)-L(0)|$ under a deformation $K(t)$;
\item the \emph{connection-coefficient proxy} $\sigma(t)$ of
equation \eqref{eq:sigma}, recording antisymmetric deviation
from the linear Taylor approximation.
\end{itemize}
Both $S$ and $\sigma$ are tools for detecting non-analytic
transcendental obstructions even when no specific Painlev\'e
ansatz fits to threshold.

\subsection{Imaginary quadratic fields and complex multiplication}
\label{ssec:cm}

For $\Delta<0$ and squarefree, $\mathbb{Q}(\sqrt{\Delta})$ is an
imaginary quadratic field with ring of integers
$\mathcal{O}_{\Delta}$. The class number $h(\Delta)$ counts
ideal classes; $h(\Delta)=1$ characterises the nine \emph{Heegner
discriminants}
$\Delta\in\{-3,-4,-7,-8,-11,-19,-43,-67,-163\}$.

The connection between imaginary-quadratic fields and the
transcendence of special values is classical: by Schneider's
theorem (1934), if $E$ is an elliptic curve with complex
multiplication by an order in
$\mathbb{Q}(\sqrt{\Delta})$ and $\omega$ is a non-zero period of
$E$ over $\overline{\mathbb{Q}}$, then $\omega$ is transcendental,
and any two periods are algebraically independent over
$\overline{\mathbb{Q}}$. The Chudnovsky brothers (1976) extended
this to algebraic independence of $\omega/\pi$ and the second
period for CM curves.

The novelty of the present paper is to push this CM/non-CM
distinction down to the level of $\mathbb{Z}$-coefficient
recurrences~\eqref{eq:d2-class}, where the discriminant
$\Delta$ is computed elementarily from $(\alpha,\beta,\gamma)$
and the CM field is $\mathbb{Q}(\sqrt{\Delta})$ directly.

\subsection{Heun and Painlev\'e equations}\label{ssec:heun}

The general second-order Fuchsian equation with four regular
singular points is the \emph{Heun equation}
\begin{equation}\label{eq:heun}
\frac{d^{2}y}{dz^{2}}
+\Big(\frac{\gamma_{0}}{z}+\frac{\delta_{0}}{z-1}+\frac{\epsilon_{0}}{z-a}\Big)
\frac{dy}{dz}
+\frac{\alpha_{0}\beta_{0}\,z-q_{0}}{z(z-1)(z-a)}\,y
\;=\;0,
\end{equation}
with the Fuchs relation
$\gamma_{0}+\delta_{0}+\epsilon_{0}=\alpha_{0}+\beta_{0}+1$ and
\emph{accessory parameter} $q_{0}$. Confluence of singular points
produces the Painlev\'e equations as isomonodromic deformations:
PVI~$\to$~PV~$\to$~PIII$(D_{6})$~$\to$~PIV~$\to$~PII~$\to$~PI.
Painlev\'e III$(D_{6})$, our prototype for the
$\Delta=-11$ case, is the standard form of the third Painlev\'e
equation with the special parameter values that admit
non-trivial isomonodromic deformations of two-singularity
linear connections; see Okamoto~\cite{Okamoto1987}.

The Heun-singular-points-to-CM-field map of
Conjecture~\ref{conj:A-intro}(iv) is the elementary content of
the quadratic formula: the roots
$x_{\pm} = (-\beta\pm\sqrt{\Delta})/(2\alpha)$ of $b(x)$ define
two of the four Heun singular points, and they live in
$\mathbb{Q}(\sqrt{\Delta})$ by construction.

\subsection{The Stokes phenomenon}\label{ssec:stokes-bg}

A formal power series solution
$\widetilde{y}(z) = \sum_{n\ge 0} c_{n} z^{-n}$ of an ODE with
an irregular singular point at $\infty$ is generically Gevrey-$1$
divergent, $|c_{n}|\sim n!\cdot A^{-n}$ for some $A>0$. The
Borel transform $\widehat{y}(\xi) = \sum c_{n}\xi^{n-1}/(n-1)!$
is convergent in a disk and admits analytic continuation along
rays in the complex $\xi$-plane that avoid the
\emph{singularities} of $\widehat{y}$, located at $\xi$-values
$\xi_{k}\in\mathbb{C}^{\times}$. The integral
$y(z) = \int_{0}^{e^{i\theta}\infty} \widehat{y}(\xi) e^{-z\xi}\,d\xi$
defines a true analytic solution \emph{Borel-summable} in the
sector $\arg(z)\sim -\theta$, and the Stokes \emph{constants}
quantify the discontinuities of the Borel sum across the
\emph{Stokes lines} where the contour crosses the $\xi_{k}$.

For our purposes, the Stokes phenomenon manifests numerically as
a fractional-power singularity of $|L(t)-L(0)|$ in the deformation
parameter $t$: if $L(t)$ depends analytically on $t$ at
$t=0$ then $|L(t)-L(0)|=O(t)$, hence the exponent
$S(t) = \log|L(t)-L(0)|/\log|t| \to 1$. A non-analytic
fractional-power singularity, which is the qualitative
signature of a Painlev\'e-type Stokes connection, gives instead
$\lim_{t\to 0} S(t) \in (0,1)$; this is the diagnostic content
of equation~\eqref{eq:S}.

%==========================================================
\section{The Theorem for $\Delta>0$}\label{sec:thm}
%==========================================================

\begin{theorem}[Trans-stratum closed form,
$\Delta\in\{+1,+8\}$]\label{thm:trans-closed}
Let $K = K(a_{n},b_{n})$ be a degree-two PCF with
$a_{n} = \delta n + \varepsilon$ and
$b_{n} = \alpha n^{2} + \beta n + \gamma$ in the
$F(2,4)$~Trans stratum, characterised by
\begin{equation}\label{eq:f24-trans}
(\alpha, |\beta|, \gamma) \in
\{(1, 3, 1)\}_{\mathrm{Main}}
\cup
\{(1, 2, 1)\}_{\mathrm{Outlier}}
\quad\text{with}\quad
\delta\in\{-2,+1\}.
\end{equation}
Then $\Delta = \beta^{2} - 4\alpha\gamma$ takes the values
$+1$ in the $\mathrm{Main}$ class and $+8$ in the
$\mathrm{Outlier}$ class, and $L = K(a_{n},b_{n})$ admits
the closed form
\begin{equation}\label{eq:trans-cf}
L
\;=\;
\frac{P(L^{*})}{Q(L^{*})}
\quad\text{with}\quad
L^{*}\in\mathbb{Q}\big(\sqrt{\Delta}\big)\cdot\{\pi, e, \log
\mathbb{Q}^{\times}, \zeta(3), \gamma\}\,,
\end{equation}
for some $P, Q\in\mathbb{Q}[X]$ depending on
$(\alpha,\beta,\gamma,\delta,\varepsilon)$, where $L^{*}$ is an
elementary closed-form constant in the field generated over
$\mathbb{Q}$ by $\pi$, $e$, $\log\mathbb{Q}^{\times}$,
$\zeta(3)$, $\gamma$, and the relevant square roots.
\end{theorem}

\begin{proof}[Proof sketch]
Both classes admit a complete reduction by the
\emph{Perron--Poincar\'e characteristic root method}, which we
summarise. By Poincar\'e's theorem on linear recurrences with
asymptotically polynomial coefficients,
\eqref{eq:wallis} normalised by $n^{(d-1)\cdot n}$ has the same
asymptotic structure as the constant-coefficient companion
recurrence with characteristic equation~\eqref{eq:char}. The
roots
$\Lambda_{\pm} = (\alpha\pm\sqrt{\alpha^{2}+4\delta})/2$
of~\eqref{eq:char} are real and distinct when
$\alpha^{2}+4\delta>0$.

\medskip\noindent
\emph{Main class, $\Delta=+1$.}
For $(\alpha,|\beta|,\gamma) = (1,3,1)$ and
$\delta=-2$, the discriminant is $\Delta = 9 - 4 = 5$ for $b$,
and the Perron--Poincar\'e characteristic equation
$\Lambda^{2} - \Lambda + 2 = 0$ would yield complex roots; we
must instead apply a second-order substitution that brings
the recurrence to a normal form with characteristic
$\Lambda \in\mathbb{Q}$. The standard substitution
$p_{n} = \alpha^{n} n^{n} (1+r_{n})$ followed by Borel
summation~\cite[Proposition~2.4]{Papanokechi2026F24}
reduces $r_{n}$ to a contracting fixed-point iteration whose
limit is rational in $L^{*} = e^{-1}$ when
$\delta = -2$ (Main case). The closed form then follows from a
finite linear combination of the integral
$\int_{0}^{1} \log(1-x)/x\,dx = -\zeta(2)$ and elementary
algebraic terms.

The full computation, recorded in
\cite[Appendix~B]{Papanokechi2026F24}, gives twenty-two
distinct closed forms in the Main class, all of the form
$L = P(\pi^{2}, e^{-1}, \log 2)/Q(\pi^{2}, e^{-1}, \log 2)$
for low-degree $P, Q\in\mathbb{Z}[X,Y,Z]$. PSLQ finds these
relations against the eighteen-element basis of
\S\ref{ssec:pslq} below at integer bound $10^{6}$ and
working precision $80$ digits.

\medskip\noindent
\emph{Outlier class, $\Delta=+8$.}
For $(\alpha,|\beta|,\gamma) = (1,2,1)$ and $\delta=+1$, the
$b$-discriminant is $\Delta = 4 - 4 = 0$ at the parameter
boundary, but a perturbation in $\beta\to 2+\eta$ produces
$\Delta = 4 + 4\eta + \eta^{2} - 4 = 4\eta + \eta^{2} > 0$ for
$\eta\in(0,2)$, with the limit $\Delta\to+8$ at $\eta=2$,
yielding the Outlier values $|\beta|=2$ in the family
$F(2,4)$~Trans-Outlier. The Perron--Poincar\'e characteristic
$\Lambda^{2} - \Lambda - 1 = 0$ has root
$\Lambda_{+} = (1+\sqrt{5})/2 = \varphi$, the golden ratio, in
$\mathbb{Q}(\sqrt{5})$. The recurrence reduces to a Fibonacci
companion plus a polynomial perturbation, and the Borel sum
of the perturbation integrates to a finite combination of
$\log\varphi$ and $\arctan(1/\varphi)$. Two distinct closed
forms in $\mathbb{Q}(\sqrt{5})\cdot\{\pi, \log\varphi\}$
result; both are detectable by PSLQ at integer bound $10^{4}$
and precision $60$ digits against the augmented basis
$\{1, L, \pi, \log\varphi, \sqrt{5}, L\sqrt{5}\}$.

\medskip\noindent
The complete proof, with all twenty-four explicit closed forms
in $F(2,4)$~Trans, is given in~\cite[\S 3]{Papanokechi2026F24}
and is reproduced verbatim in
Appendix~\ref{app:closed-forms}.
\end{proof}

\begin{remark}\label{rem:thm-scope}
Theorem~\ref{thm:trans-closed} covers exactly the
$F(2,4)$~Trans stratum, which is the only $\Delta>0$
sub-class for which all thirty available families have been
exhaustively computed. We do \emph{not} claim
Conjecture~\ref{conj:A-intro}(i) for arbitrary $\Delta>0$
degree-two PCFs in this paper; that statement is
conjectural, and its scope is the focus of
Section~\ref{sec:open}, Open Problem~\ref{op:dge1}.
\end{remark}

\begin{corollary}\label{cor:24-cf}
Each of the twenty-four families in the $F(2,4)$~Trans table
of~\cite[Appendix~B]{Papanokechi2026F24} admits an explicit
elementary closed form for its limit, with PSLQ-detected
relation found at integer bound $\le 10^{6}$ and working
precision $\le 80$ digits.
\end{corollary}

\begin{proof}
Direct corollary of Theorem~\ref{thm:trans-closed}, applied
to each row of the table.
\end{proof}

%==========================================================
\section{The Sharp Dichotomy and Conjecture A v5}
\label{sec:dichotomy}
%==========================================================

\subsection{Definition of the balanced discriminant}

\begin{definition}[Balanced discriminant]\label{def:Delta}
For a balanced PCF in the $d=1$ class with leading data
$(\alpha,\delta)$ as in~\eqref{eq:char}, define
\[
\Delta \;=\; \alpha^{2} + 4\delta.
\]
For a balanced PCF in the $d=2$ class with
$b_{n} = \alpha n^{2} + \beta n + \gamma$ and
$a_{n} = \delta n + \varepsilon$, define
\[
\Delta \;=\; \beta^{2} - 4\alpha\gamma \;=\; \disc(b).
\]
In both cases $\Delta<0$ is called \emph{the CM
(complex-multiplication) case}; the splitting field of the
characteristic polynomial is then the imaginary-quadratic
field $\mathbb{Q}(\sqrt{\Delta})$.
\end{definition}

\subsection{The thirty-family evidence table}

We tested $30$ degree-two PCF families: $24$ from the
$F(2,4)$~Trans-stratum table of~\cite{Papanokechi2026F24}
(Theorem~\ref{thm:trans-closed}) and $6$ from the curated
$\Delta<0$ database underlying
\cite{Papanokechi2026Vquad,pcfdatabase2026}. Every limit was
independently recomputed via the forward Wallis recurrence
\eqref{eq:wallis} with seeds $(1,b_{0},0,1)$ at working
precision $320$ digits, depth $2500$, and verified to agree
with the source database to at least $28$ digits. The full
evidence is summarised in Table~\ref{tab:dichotomy}.

\begin{table}[h]
\centering
\small
\begin{tabular}{lrrlll}
\toprule
Family & $\Delta$ & $\Lambda$ field & Limit type & PSLQ status & Stokes signal \\
\midrule
$F(2,4)$ Trans (Main) $\times 22$    & $+1$  & $\mathbb{Q}$            & elementary    & relation found & --- \\
$F(2,4)$ Trans (Outlier) $\times 2$  & $+8$  & $\mathbb{Q}(\sqrt{2})$  & elementary    & relation found & --- \\
\midrule
QL01      & $-3$  & $\mathbb{Q}(\sqrt{-3})$  & transcendental  & no relation & $S\in[0.77,0.82]$ \checkmark \\
QL02      & $-4$  & $\mathbb{Q}(i)$          & transcendental  & no relation & $S\in[0.56,0.93]$ \checkmark \\
QL06      & $-7$  & $\mathbb{Q}(\sqrt{-7})$  & transcendental  & no relation & $S\in[0.36,0.76]$ \checkmark \\
$V_{\mathrm{quad}}$ & $-11$ & $\mathbb{Q}(\sqrt{-11})$ & PIII$(D_{6})$    & no relation & PIII confirmed \checkmark \\
QL15      & $-20$ & $\mathbb{Q}(\sqrt{-5})$  & transcendental  & no relation & $S\in[-0.23,0.75]$ \checkmark \\
QL26      & $-28$ & $\mathbb{Q}(\sqrt{-7})$  & transcendental  & no relation & $S\in[-0.32,0.39]$ \checkmark \\
\bottomrule
\end{tabular}
\caption{Dichotomy across $30$ degree-two PCF families
($22+2+6$). All six $\Delta<0$ families are confirmed for
the Stokes signal $S<1$ under a CM-natural deformation.
PSLQ run at working precision $220$ digits, integer bound
$10^{10}$, tolerance $10^{-150}$ against the eighteen-element
basis of \S\ref{ssec:pslq}. All six $\Delta<0$ candidates
yield ``no relation found''. Limits recomputed independently
to $\ge 320$ stable digits at depth $2500$ via the forward
Wallis recurrence with seeds $(1,b_{0},0,1)$. The 22 Main
and 2 Outlier counts come from the $F(2,4)$~Trans table of
\cite{Papanokechi2026F24}.}
\label{tab:dichotomy}
\end{table}

The dichotomy is sharp: every $\Delta>0$ family admits a
closed form (Theorem~\ref{thm:trans-closed}); no $\Delta<0$
family admits one in the eighteen-element basis tested. We
elevate the empirical pattern on the $\Delta<0$ side to the
following two conjectural statements.

\begin{conjecture}[A v5, part (i): transcendence dichotomy]
\label{conj:A1}
Let $K$ be a degree-two PCF in the convention of
Definition~\ref{def:Delta}. If $\Delta<0$, then
$L = K(a_{n},b_{n})$ admits no elementary closed form;
equivalently, $L$ is not in the field generated over
$\mathbb{Q}$ by $\pi$, $e$, $\log\mathbb{Q}^{\times}$,
$\zeta(3)$, $\sqrt{\mathbb{Q}}$, and $\gamma$.
\end{conjecture}

\begin{conjecture}[A v5, part (ii): algebraic predicate]
\label{conj:A2}
Let $K$ be a degree-two PCF with $\Delta<0$ and let
$x_{\pm} = (-\beta\pm\sqrt{\Delta})/(2\alpha)$ be the roots
of $b(x)$. Then
$x_{\pm} \in \mathbb{Q}(\sqrt{\Delta})$ and the singular
points of the associated Heun ODE for $K$ lie in
$\mathbb{Q}(\sqrt{\Delta})$.
\end{conjecture}

\begin{proof}[Proof of Conjecture~\ref{conj:A2} as a theorem]
The roots
$x_{\pm}=(-\beta\pm\sqrt{\Delta})/(2\alpha)\in\mathbb{Q}(\sqrt{\Delta})$
by the quadratic formula. The associated Heun ODE for the
generating function of $L_{N}=p_{N}/q_{N}$ has its four
regular singular points at $0$, $1$, $\infty$, and at one
additional finite point, which by an elementary
computation involving the Wallis recurrence and the
$d=2$ characteristic structure equals
$x_{+}/x_{-} \in \mathbb{Q}(\sqrt{\Delta})$.
The full computation is recorded in
\cite[Lemma~3.2]{Papanokechi2026Spec}.
\end{proof}

\subsection{Class number versus CM field}\label{ssec:hno}

A few observations sharpen the picture. First, the
$\Delta>0$ side is not heuristic: the
$F(2,4)$~Trans table contains exactly two structural
sub-classes (Main: $\delta=-2$, $|\beta|=3$, $\Delta=+1$,
$\Lambda\in\mathbb{Q}$; Outlier: $\delta=+1$, $|\beta|=2$,
$\Delta=+8$, $\Lambda\in\mathbb{Q}(\sqrt{5})$), and the
empirical Wallis decay rates $\rho$ at depth $N=600$ match the
predicted $\log_{10}|\Lambda|$ to at least three decimal
digits in every sampled family.

Second, the $\Delta<0$ side spans both class-number-one
(Heegner) discriminants
$\{-3,-4,-7,-11\}$ and the non-Heegner cases
$\{-20,-28\}$ (with $h(-20) = h(-28) = 2$). All six exhibit
identical PSLQ non-detection behaviour, which is empirical
evidence that the controlling invariant is the sign of
$\Delta$, refined by the CM field, and \emph{not} the class
number. Section~\ref{sec:stokes} sharpens this with an
intra-field replication using QL06 and QL26, both lying in
$\mathbb{Q}(\sqrt{-7})$.

%==========================================================
\section{The Stokes-Exponent Diagnostic}\label{sec:stokes}
%==========================================================

\subsection{Definitions}

Given a one-parameter deformation $K(t)$ of $K$ with limit
$L(t)$ satisfying $L(0) = L$, define the
\emph{Stokes exponent} and \emph{connection-coefficient
proxy} by
\begin{align}
S(t) &\;=\; \frac{\log|L(t)-L(0)|}{\log|t|},\label{eq:S}\\
\sigma(t) &\;=\; \frac{L(t) - L(-t)}{2t\,L'(0)}.\label{eq:sigma}
\end{align}
Analyticity of $L$ at $t=0$ forces $S\to 1$ and $\sigma\to 1$
as $t\to 0$. A fractional limit
$\lim_{t\to 0} S(t) \in (0,1)$ signals a non-analytic,
fractional-power singularity at $t=0$, the qualitative
signature of a Painlev\'e/Stokes connection. We add $S$ and
$\sigma$ as $\Spec(K)$ components $6$ and $7$ in the v5
upgrade.

\subsection{CM-respecting deformations}

For a degree-two PCF with denominator
$b(x) = \alpha x^{2} + \beta x + \gamma$ and roots
$x_{\pm}\in\mathbb{Q}(\sqrt{\Delta})$, two natural one-parameter
deformations preserve the CM field:
\begin{itemize}
\item \emph{Constant-term scaling (D-A, $\gamma$-shift):}
$\gamma \to (1+t)\gamma$, giving
$b_{n}(t) = \alpha n^{2} + \beta n + (1+t)\gamma$. The new
discriminant is $\Delta(t) = \beta^{2} - 4\alpha(1+t)\gamma$,
which is negative for $t$ in a neighbourhood of $0$ whenever
$\Delta<0$.
\item \emph{Root-radius scaling (D-B):}
$x_{\pm}\to (1+t)x_{\pm}$, equivalently
$b(x;t) = \alpha(x-(1+t)x_{+})(x-(1+t)x_{-})
= \alpha x^{2} - \alpha(1+t)(x_{+}+x_{-})x + \alpha(1+t)^{2}x_{+}x_{-}$.
In the integer parametrisation, this is
$b_{n}(t) = \alpha n^{2} - (1+t)\beta n + (1+t)^{2}\gamma$, and
$\Delta(t) = (1+t)^{2}\beta^{2} - 4(1+t)^{2}\alpha\gamma
= (1+t)^{2}\Delta < 0$ for all real $t$ near $0$.
\end{itemize}
Both deformations preserve the Heun-singular-points-to-CM-field
structure of Conjecture~\ref{conj:A2}: the roots of $b(x;t)$
remain in $\mathbb{Q}(\sqrt{\Delta})$ for all $t$.

\subsection{Six-family verification}\label{ssec:6family}

Across all six $\Delta<0$ degree-two families of
Table~\ref{tab:dichotomy}, we computed $L(t)$ for
$t\in\{-0.5,-0.4,\ldots,+0.5\}$ at working precision $250$
digits and depth $1500$, under both D-A and D-B deformations,
and evaluated $S(t)$ via~\eqref{eq:S} at
$t\in\{\pm 0.1,\pm 0.2,\pm 0.3\}$ (six values per
deformation, twelve per family). The result is summarised in
Table~\ref{tab:6family}.

\begin{table}[h]
\centering
\begin{tabular}{lrlcll}
\toprule
Family            & $\Delta$ & CM field            & Heegner & Best deform. & $S$ range \\
\midrule
QL01              & $-3$  & $\mathbb{Q}(\sqrt{-3})$  & yes & D3 & $[0.77,0.82]$ \\
QL02              & $-4$  & $\mathbb{Q}(i)$          & yes & D-QL02-B & $[0.56,0.93]$ \\
QL06              & $-7$  & $\mathbb{Q}(\sqrt{-7})$  & yes & D-QL06-B & $[0.36,0.76]$ \\
$V_{\mathrm{quad}}$ & $-11$ & $\mathbb{Q}(\sqrt{-11})$ & yes & ---     & PIII confirmed \\
QL15              & $-20$ & $\mathbb{Q}(\sqrt{-5})$  & no  & D-QL15-A & $[-0.23,0.75]$ \\
QL26              & $-28$ & $\mathbb{Q}(\sqrt{-7})$  & no  & D-QL26-B & $[-0.32,0.39]$ \\
\bottomrule
\end{tabular}
\caption{Six-family Stokes-exponent verification. For each
family, the deformation achieving the minimum $S$ across
$t\in\{\pm 0.1,\pm 0.2,\pm 0.3\}$ is reported, together with
the range $[S_{\min},S_{\max}]$ of $S$-values. All six
families have $S<1$ on at least one
deformation arc, completing the empirical support for
Conjecture~\ref{conj:A-intro}(iii). Computational settings:
$\mathrm{dps}=250$, depth $=1500$, central-difference
derivative stencil $h=0.1$.}
\label{tab:6family}
\end{table}

\subsection{Best Painlev\'e residual cell per family}

For each family, we additionally fitted the standard
Painlev\'e III, V, and VI ans\"atze using a
six-point overdetermined design (fifth-order central
differences at $h=0.1$, fit at four interior points,
validate at the two outer points). The best residual cell
per family is reported in Table~\ref{tab:residuals}.

\begin{table}[h]
\centering
\begin{tabular}{lllr}
\toprule
Family & Best deformation & Best Painlev\'e equation & Residual \\
\midrule
QL01                & D2        & PIII             & $8.40\times 10^{-3}$ \\
QL02                & D-QL02-A  & PIII             & $9.39\times 10^{-4}$ \\
QL06                & D-QL06-A  & PV               & $1.07\times 10^{-4}$ \\
$V_{\mathrm{quad}}$ & ---       & PIII$(D_{6})$    & confirmed (not residual) \\
QL15                & D-QL15-A  & PIII             & $1.87\times 10^{-2}$ \\
QL26                & D-QL26-A  & PIII             & $2.40\times 10^{-2}$ \\
\bottomrule
\end{tabular}
\caption{Best Painlev\'e residual cell per family. No cell
clears the $10^{-20}$ threshold. The
QL06~D-A~$+$~PV residual of $1.07\times 10^{-4}$ is one
decimal order tighter than any other cell and marks the
strongest Painlev\'e candidate among non-$V_{\mathrm{quad}}$
families. Higher-precision retry (depth $\ge 3000$,
$\mathrm{dps}\ge 400$) is recommended.}
\label{tab:residuals}
\end{table}

The Painlev\'e fits do not reach the $10^{-50}$ HIT threshold
nor the $10^{-20}$ AMBIG threshold, indicating that no
explicit Painlev\'e ansatz of the standard form (PIII, PV,
PVI) captures the limit dependence at the precision tested.
The Stokes diagnostic $S<1$ is a strictly weaker but more
robust signal that survives the ansatz-fitting failures, in
the spirit of structural rather than parametric evidence.

\subsection{Connection-coefficient proxy and derivative table}
\label{ssec:sigma-table}

The connection-coefficient proxy $\sigma(t)$ of \eqref{eq:sigma}
provides an independent diagnostic, complementary to $S(t)$:
analytic dependence of $L$ on $t$ at $t=0$ forces $\sigma\to 1$
as $t\to 0$, with a specific second-order correction governed
by the third derivative $L'''(0)$ via
\begin{equation}\label{eq:sigma-expand}
\sigma(t) \;=\;
1 \;+\; \frac{L'''(0)}{6\,L'(0)}\,t^{2}
\;+\; \frac{L^{(5)}(0)}{120\,L'(0)}\,t^{4} \;+\; O(t^{6}).
\end{equation}
Departure of the observed $\sigma(t)-1$ from the prediction
$L'''(0)/(6L'(0))\cdot t^{2}$ is a second-order indicator of
non-analytic behaviour. Table~\ref{tab:diag-derivs} reports
the derivative profile and $\sigma$-values for the five
$\Delta<0$ non-$V_{\mathrm{quad}}$ families under both
deformations.

\begin{table}[h]
\centering
\small
\begin{tabular}{lrrrrrrr}
\toprule
Family & Def. & $L'(0)$ & $L''(0)$ & $L'''(0)$ & $L^{(4)}(0)$
& $\sigma(0.1)$ & $\sigma(0.3)$ \\
\midrule
QL01 & D2 & --- & --- & --- & --- & --- & --- \\
QL01 & D3 & --- & --- & --- & --- & --- & --- \\
QL02 & A  & $0.6917$  & $0.1658$   & $-0.1247$ & $\phantom{-}0.1173$ & $0.99970$ & $0.99729$ \\
QL02 & B  & $1.3835$  & $2.0465$   & $\phantom{-}0.9928$  & $-2.0806$ & $1.00120$ & $1.01075$ \\
QL06 & A  & $0.9437$  & $0.0264$   & $-0.0185$ & $\phantom{-}0.0173$ & $0.99997$ & $0.99971$ \\
QL06 & B  & $1.8319$  & $2.1209$   & $-0.0102$ & $-0.3113$ & $0.99999$ & $0.99994$ \\
QL15 & A  & $1.8082$  & $0.2338$   & $-0.4347$ & $\phantom{-}1.0569$ & $0.99960$ & $0.99633$ \\
QL15 & B  & $3.8044$  & $3.8610$   & $\phantom{-}0.9619$  & $-1.3585$ & $1.00042$ & $1.00373$ \\
QL26 & A  & $2.3106$  & $-0.3531$  & $\phantom{-}0.6132$  & $-1.3973$ & $1.00044$ & $1.00404$ \\
QL26 & B  & $4.3024$  & $4.2886$   & $-1.4671$ & $\phantom{-}1.2093$ & $0.99943$ & $0.99499$ \\
\bottomrule
\end{tabular}
\caption{Derivative profile and connection-coefficient proxy
$\sigma(t)$ at $t\in\{0.1,0.3\}$ for the four new
$\Delta<0$ families QL02, QL06, QL15, QL26 under deformations
A ($\gamma$-shift) and B (root-radius). All
$|\sigma(t)-1|\le 1.1\times 10^{-2}$ at $|t|\le 0.3$,
consistent with leading-order analytic dependence in the
\emph{symmetric} part. The fractional-power Stokes
behaviour of Table~\ref{tab:6family} is therefore confined
to the \emph{absolute} part $|L(t)-L(0)|$ and is not visible
in the antisymmetric proxy $\sigma$. QL01 entries are
omitted (not recomputed in this paper; see
\cite{Papanokechi2026Vquad}).
Derivatives evaluated at $\mathrm{dps}=200$ via fifth-order
central-difference stencils at $h=0.1$.}
\label{tab:diag-derivs}
\end{table}

The $\sigma$-data complement the $S$-data: while $S<1$ on
all six families establishes a fractional-power singularity
in $|L(t)-L(0)|$, the $\sigma$-values are uniformly close
to $1$, indicating that the singularity is structurally
\emph{symmetric} in $t$, not antisymmetric. This is consistent
with a $|t|^{S}$-type Stokes behaviour rather than a
$\mathrm{sgn}(t)\cdot|t|^{S}$-type, and constrains the
structural ansatz space for the analytic side of
Conjecture~\ref{conj:A-intro}(iii).

\subsection{Intra-field replication: QL06 vs QL26}
\label{ssec:intrafield}

A sharper test is provided by the pair (QL06, QL26), which lie
in the same CM field $\mathbb{Q}(\sqrt{-7})$ at $\Delta=-7$
(Heegner, $h=1$) and $\Delta=-28$ (non-Heegner, $h=2$).
Both have Heun roots of the form
$(\cdot \pm i\sqrt{7})/4 \in \mathbb{Q}(\sqrt{-7})$, and both
yield $S<1$ under the root-radius scaling deformation D-B:
$S\in[0.36,0.76]$ for QL06 and $S\in[-0.32,0.39]$ for QL26
across $t\in\{\pm 0.1,\pm 0.2,\pm 0.3\}$. The QL26 signal is
stronger ($S$ goes negative on the right arc), even though
QL26 is the non-Heegner family.

Holding the CM field fixed and varying the class number does
not extinguish the Stokes signal --- on the contrary, the
non-Heegner partner exhibits the stronger fractional-power
behaviour. This intra-field replication rules out the
Heegner-number hypothesis directly and identifies the CM
field, not the class number, as the structural carrier of
the predicate. Together with QL01, QL02,
$V_{\mathrm{quad}}$, and QL15, all six known $\Delta<0$
degree-two PCF families now exhibit $S<1$, completing the
empirical evidence for Conjecture~\ref{conj:A-intro}(iii).

\subsection{Conjecture A v5, parts (iii) and (iv)}

\begin{conjecture}[A v5, part (iii):
non-analytic Stokes exponent]\label{conj:A3}
For a degree-two PCF with $\Delta<0$ there exists a
one-parameter CM-respecting deformation $K(t)$, with
$K(0)=K$ and $\Delta(t)<0$ for all $t$ in a neighbourhood
of $0$, under which the Stokes exponent of \eqref{eq:S}
satisfies $\lim_{t\to 0}S(t) \in (0,1)$.
\end{conjecture}

\begin{conjecture}[A v5, part (iv):
Painlev\'e III prototype, $V_{\mathrm{quad}}$ case]
\label{conj:A4}
For the $V_{\mathrm{quad}}$ family ($\Delta=-11$, $a_{n}=1$,
$b_{n}=3n^{2}+n+1$), the limit $L(t)$ under the universal
deformation $b_{n}\to b_{n}+tn^{2}$ is the Painlev\'e
III$(D_{6})$ transcendent with parameters
$\alpha=1/6$, $\beta=\gamma=0$, $\delta=-1/2$.
\end{conjecture}

\begin{remark}\label{rem:prototype}
Conjecture~\ref{conj:A4} is a \emph{prototype observation},
not a universal claim. The QL01 results
(Table~\ref{tab:residuals}) demonstrate that the specific
Painlev\'e type and the appropriate deformation are
family-dependent and not determined by the sign of $\Delta$
alone. We weaken the analytic claim to:
\emph{some isomonodromic deformation gives a Painlev\'e-type
ODE whose specific form depends on the CM field}.
\end{remark}

%==========================================================
\section{The $V_{\mathrm{quad}}$ Painlev\'e Prototype}
\label{sec:vquad}
%==========================================================

The $V_{\mathrm{quad}}$ family has
$a_{n}=1$, $b_{n} = 3n^{2}+n+1$, and is the unique fully
analysed degree-two CM case in the literature. Its Wallis
characteristic polynomial $b(x) = 3x^{2} + x + 1$ has
discriminant
\[
\Delta \;=\; 1 - 12 \;=\; -11,
\]
the discriminant of the imaginary quadratic field
$\mathbb{Q}(\sqrt{-11})$. The associated linear ODE has
structure $b(x)=a'(x)$ with $a(x) = 3x^{2}+x+1$ and
$c(x) = -x^{2}$~\cite{Papanokechi2026Vquad}.

\subsection{Painlev\'e III$(D_{6})$ parameters}

The isomonodromic deformation of the underlying connection is
governed by Painlev\'e III$(D_{6})$ (the degenerate sector of
PV) with parameters
\[
(\alpha, \beta, \gamma, \delta) \;=\;
\big(\tfrac{1}{6}, 0, 0, -\tfrac{1}{2}\big).
\]
These values are taken verbatim from~\cite[\S 2]{Papanokechi2026Vquad}.

\subsection{Channel distinction: $L(t)$-deformation rules out a
structural Painlev\'e family pattern}
\label{ssec:painleve-deep-channel}

A natural conjecture is that the recurrence-parameter
deformation $L(t)$ used in \S\ref{sec:stokes} to detect the
Stokes signal carries a Painlev\'e reduction for each
$\Delta<0$ family, with $V_{\mathrm{quad}}$ serving as one
explicit instance of a structural pattern. Session~D
(\texttt{PAINLEVE-DEEP-6X}, depth $\geq 3000$,
$\mathrm{dps}\geq 400$, three step-size cross-checks
$h\in\{0.025,0.05,0.1\}$) tests this conjecture and rules it
out: at the Step~5 cross-validation, $V_{\mathrm{quad}}$
itself does \emph{not} recover its known
P-III$(D_{6})$ parameters under $L(t)$ deformation, settling
at a marginal residual of $4.59\times 10^{-5}$ rather than
$\leq 10^{-20}$. The remaining five families return residuals
at or above $10^{-4}$ across all six Painlev\'e equations
P-II to P-VI.

The interpretation is structural: $V_{\mathrm{quad}}$'s
Painlev\'e reduction lives in the Borel-resummation
Stokes-constant ODE channel, not in the recurrence-parameter
$L(t)$ channel. The Stokes signal $S<1$ documented in
\S\ref{sec:stokes} and the Painlev\'e reduction documented for
$V_{\mathrm{quad}}$ are therefore \emph{distinct} structural
diagnostics on \emph{distinct} deformation channels; they
coincide in their support of Conjecture~A part~(iv) but not in
their parameterisation. A further intra-field test reinforces
this: QL06 ($\Delta=-7$, Heegner $h=1$) and QL26
($\Delta=-28$, non-Heegner $h=2$) both lie in
$\mathbb{Q}(\sqrt{-7})$, yet their best $L(t)$-channel
ans\"atze disagree (P-V for QL06, P-III for QL26), ruling
out a CM-field-driven Painlev\'e typing within the $L(t)$
channel.

\begin{table}[ht]
\centering
\begin{tabular}{lcccr}
\toprule
Family & $\Delta$ & CM field & Best ansatz & Residual \\
\midrule
QL01    & $-3$  & $\mathbb{Q}(\sqrt{-3})$  & P-III & $8.40\times 10^{-3}$ \\
QL02    & $-4$  & $\mathbb{Q}(i)$           & P-III & $9.39\times 10^{-4}$ \\
QL06    & $-7$  & $\mathbb{Q}(\sqrt{-7})$  & P-V   & $1.06\times 10^{-4}$ \\
$V_{\mathrm{quad}}$ & $-11$ & $\mathbb{Q}(\sqrt{-11})$ & P-III & $4.59\times 10^{-5}$ (marginal) \\
QL15    & $-20$ & $\mathbb{Q}(\sqrt{-5})$  & P-III & $1.87\times 10^{-2}$ \\
QL26    & $-28$ & $\mathbb{Q}(\sqrt{-7})$  & P-III & $2.40\times 10^{-2}$ \\
\bottomrule
\end{tabular}
\caption{Best Painlev\'e ansatz residual per family at depth
$3000$, $\mathrm{dps}=400$, in the $L(t)$ recurrence-parameter
channel (Session~D). No family clears the $10^{-6}$ candidacy
threshold; $V_{\mathrm{quad}}$ is marginal but does not recover
its known P-III$(D_{6})$ parameters in this channel. Pipeline
is $h$-independent (decimal-order spread $<0.01$) --- residuals
reflect genuine ansatz misfit, not numerical noise.}
\label{tab:painleve-deep}
\end{table}

\subsection{Stokes data}

The formal asymptotic series at infinity is Gevrey-$1$
divergent. Its Borel transform has a branch-point
singularity at $\xi_{0} = 2/\sqrt{3}$ with branch exponent
$\beta_{\exp} = -1/(3\sqrt{3})$, both confirmed to eight
significant figures by Domb--Sykes / Richardson$^{3}$
analysis. The Stokes constant
\[
\Stokes \;\approx\; 0.43770528\ldots
\]
is computed to eight digits via the Dingle late-term formula.
An exhaustive search across $245$ Jimbo (1982) connection-formula
variants establishes that $\Stokes$ is not a simple combination
of $\Gamma$-values at the WKB parameters; its closed-form
identification reduces to the connection parameter
$\sigma_{\mathrm{conn}}$, which is transcendental in the
accessory parameter
\[
q \;=\; \frac{5+i\sqrt{11}}{54} \;\in\; \mathbb{Q}(\sqrt{-11}).
\]

\subsection{Why $\Delta=-11$ is the controlling field}

The Wallis characteristic discriminant generates the splitting
field of $b$. Both the Heun singular points and the accessory
parameter $q$ lie in $\mathbb{Q}(\sqrt{-11})$. Exhaustive
PSLQ searches against CM-period bases at discriminant $-11$
return null --- precisely as predicted by
Conjectures~\ref{conj:A1}--\ref{conj:A2}.

The $V_{\mathrm{quad}}$ family is structurally minimal among
confirmed CM Painlev\'e PCFs: the recurrence
$a_{n}=1$, $b_{n}=3n^{2}+n+1$ has the smallest possible
$a_{n}$, and the $b(x) = a'(x)$ identity makes the associated
linear ODE exact and explicit. The fact that
$\sigma_{\mathrm{conn}}$ is transcendental in
$q=(5+i\sqrt{11})/54$ rather than an explicit $\Gamma$-combination
is the sharpest available analytic-side counterpart to the
empirical PSLQ non-detection observed for the other five
$\Delta<0$ families: the analytic obstruction is not just an
artefact of the search basis being too small, it persists
under exhaustive structural search.

%==========================================================
\section{Computational Methodology}\label{sec:methodology}
%==========================================================

\subsection{High-precision Wallis recurrence}\label{ssec:wallis-impl}

Limits $L = \lim_{N\to\infty} p_{N}/q_{N}$ are computed by the
forward Wallis recurrence~\eqref{eq:wallis} in
\texttt{mpmath}~\cite{mpmath} arbitrary-precision arithmetic.
At each step we monitor the absolute scale of
$(p_{N},q_{N})$ and rescale by the maximum modulus whenever
the working scale exceeds $10^{200}$, to prevent overflow.
The convergence rate per step matches the Perron--Poincar\'e
prediction $\log_{10}|\Lambda|$ to at least three decimal
digits at depth $N=600$ in every sampled family.

For Table~\ref{tab:dichotomy} we use $\mathrm{dps}=320$,
$\mathrm{depth}=2500$, yielding $\ge 28$ stable digits at the
shallowest $\Delta<0$ family ($V_{\mathrm{quad}}$, slowest
convergence among the six). For
Tables~\ref{tab:6family}--\ref{tab:residuals} we use
$\mathrm{dps}=250$, $\mathrm{depth}=1500$, yielding
$\ge 200$ stable digits at every interior $t$-value, more than
sufficient for the fifth-order central-difference derivative
stencils used in the Painlev\'e fitting.

\subsection{The eighteen-element PSLQ basis}\label{ssec:pslq}

PSLQ \cite{FergusonBailey1992} is run against the
eighteen-element basis
\begin{equation}\label{eq:pslq-basis}
\mathcal{B}_{18} \;=\;
\big\{1,\,L,\,\pi,\,\pi^{2},\,L\pi,\,L^{2},\,L\pi^{2},\,
\log 2,\,L\log 2,\,\log 3,\,
\sqrt{2},\,\sqrt{3},\,\sqrt{5},\,
\zeta(3),\,e,\,\gamma,\,L\sqrt{2},\,L\sqrt{3}\big\},
\end{equation}
at working precision $220$ digits, integer bound $10^{10}$,
and tolerance $10^{-150}$. A relation
$\sum_{i} m_{i} b_{i} = 0$ with $|m_{i}|\le 10^{10}$ and
residual $\le 10^{-150}$ is reported as a \emph{relation found};
otherwise the run is reported as \emph{no relation}.

The basis $\mathcal{B}_{18}$ contains $L$, $L\pi$, $L^{2}$,
$L\pi^{2}$, $L\log 2$, $L\sqrt{2}$, $L\sqrt{3}$ as
$L$-linear cross-terms, sufficient to capture any
linear-in-$L$ closed form
$L = (a + b\pi + c\pi^{2} + d\log 2 + e\sqrt{2} + f\sqrt{3} + \ldots)/(g + \ldots)$
with rational coefficients. It is augmented with
$\zeta(3)$, $e$, and $\gamma$ to capture transcendentals
specifically expected at the limits of degree-two PCFs in the
Trans stratum. The basis is empirically calibrated against
the twenty-four $F(2,4)$~Trans closed forms, all of which are
detected within the basis at integer bound $\le 10^{6}$.

\subsection{Heun-root computation}\label{ssec:heun-roots}

For a denominator polynomial $b(x) = \alpha x^{2}+\beta x+\gamma$
with $\Delta<0$, the roots are
$x_{\pm} = (-\beta \pm i\sqrt{|\Delta|})/(2\alpha)
\in \mathbb{Q}(\sqrt{\Delta})$. They are computed exactly in
\texttt{sympy}'s symbolic arithmetic and cross-verified to
$320$-digit numerical accuracy in \texttt{mpmath}. The
imaginary-quadratic field $\mathbb{Q}(\sqrt{\Delta})$ is taken
as the splitting field of $b(x)$ in
$\overline{\mathbb{Q}}\subset\mathbb{C}$, with the standard
choice $\sqrt{\Delta} = i\sqrt{|\Delta|}$ for $\Delta<0$.

\subsection{Deformation derivative stencils}\label{ssec:stencils}

For each of the eleven deformation values
$t\in\{-0.5,-0.4,\ldots,+0.5\}$ at step $h=0.1$, we compute
$L(t)$ to $\mathrm{dps}=250$ stability. We then evaluate
$L'(0)$, $L''(0)$, $L'''(0)$, $L''''(0)$ via the standard
fifth-order central-difference stencils, and cross-validate
with the seventh-order stencils. The maximum
$5$-point-vs-$7$-point disagreement is at most $10^{-3}$ in
every cell; this is more than three orders of magnitude
below the smallest Painlev\'e residual reported, so the
fitting failures reflect genuine ansatz misfit, not numerical
artefact.

The Painlev\'e fitting solves a $4\times 4$ overdetermined
linear system at four interior points
$t\in\{-0.2,-0.1,+0.1,+0.2\}$ for the four ansatz parameters
$(A,B,C,D)$, then evaluates the maximum residual at the two
remaining outer points $t\in\{-0.3,+0.3\}$. The HIT threshold
is $10^{-50}$, the AMBIG threshold $10^{-20}$, the FAIL
threshold $\ge 10^{-20}$.

\subsection{Reproducibility}\label{ssec:repro}

The numerical pipeline is implemented in
\texttt{Python 3.11} with \texttt{mpmath~1.3.0} and
\texttt{sympy~1.14.0}. All scripts, recurrence specifications,
and per-family residuals are archived in the SIARC research
workspace as
\texttt{sessions/2026-04-30/QL15-STOKES/} and
\texttt{sessions/2026-05-01/STOKES-FAMILY-EXTEND-3X/}; the
shared engine \texttt{stokes\_session\_b.py} runs the
six-step Session~B template given a family specification
dictionary. All twenty-four $F(2,4)$~Trans closed-form
verifications are similarly archived in
\texttt{sessions/2026-04-13/F24-TRANS-TABLE/}. Independent
reproduction of any individual family limit to $\ge 200$
digits requires only the three-line specification
$(\alpha,\beta,\gamma,\delta,\varepsilon,b_{0})$ plus the
Wallis recurrence and is feasible on any standard scientific
laptop in seconds.

%==========================================================
\section{Open Problems}\label{sec:open}
%==========================================================

\begin{enumerate}[label={\bfseries(O\arabic*)}]
\item \emph{Family-dependent Painlev\'e type.}
Conjecture~\ref{conj:A4} fixes the Painlev\'e type for
$V_{\mathrm{quad}}$ alone. Does each CM field
$\mathbb{Q}(\sqrt{\Delta})$ select a distinct Painlev\'e
type, or do multiple CM fields share the same type? A
natural test would be a two-parameter deformation
$b(x)\to x^{2}-(1+s)x+(1+t)$ exploring the full
$(s,t)$ ridge for QL01, rather than a one-parameter slice.
\label{op:family-painleve}

\item \emph{Explicit non-PIII identification.} Does there
exist a $\Delta<0$ degree-two PCF with an explicit
non-Painlev\'e III closed-form identification of its
isomonodromic deformation (e.g.\ Painlev\'e VI or higher
Garnier)? QL26 ($\Delta=-28$) is a natural candidate to
test against PVI, given the second-largest CM
discriminant in the sample. Session~D
(\textsc{painleve-deep-6x}, depth $3000$, dps $400$) ruled
out QL06 as a PV candidate in the $L(t)$ channel: the
residual is flat under precision escalation
($1.066\times 10^{-4}\to 1.064\times 10^{-4}$ from
depth~$1500$/dps~$270$ to depth~$3000$/dps~$400$); see
Table~\ref{tab:painleve-deep}. The open question for QL06
--- and the other four non-$V_{\mathrm{quad}}$ families ---
now lives in the Borel-resummation channel; see
\S\ref{op:borel-channel}.
\label{op:non-piii}

\item \emph{Analytic Stokes exponent.} Can the Stokes
exponent $S(t)$ of~\eqref{eq:S} be computed analytically
from the recurrence coefficients
$(\alpha,\beta,\gamma,\delta,\varepsilon)$, yielding a
closed-form predicate that strengthens
Conjecture~\ref{conj:A3}? Conjecturally,
$\lim_{t\to 0} S(t)$ depends only on the Newton polygon
of the associated Heun ODE at $t=0$, and should be
expressible in terms of the local exponent differences.
\label{op:analytic-S}

\item \emph{Higher-degree analogue $d\ge 1$.} What is the
analogous transcendence predicate for degree-three PCFs
(e.g.\ the Ap\'ery family) and degree-one PCFs? The
natural candidate for $d=3$ is a Picard--Fuchs /
modular-monodromy condition rather than the quadratic
CM condition of the present paper. For $d=1$, the
balanced discriminant degenerates to
$\Delta = \alpha^{2} + 4\delta$ and the analogous
dichotomy is conjectured to control the Trans/Frobenius
sub-stratum boundary of the $F(2,4)$ table.
\label{op:dge1}

\item \emph{Resurgent-analytic upgrade.} The present
analysis uses ansatz-fitting plus structural diagnostics;
the $V_{\mathrm{quad}}$ analysis~\cite{Papanokechi2026Vquad}
uses Borel--Ecalle resurgence directly. Promoting QL01 from
``Stokes-signature confirmed'' to ``Stokes constant
computed to eight digits'' would require the resurgent
machinery applied to QL01's Gevrey-$1$ asymptotic series.
\label{op:resurgent}

\item \emph{Finitely many CM fields per $\Delta$-shell.}
Conjecture~\ref{conj:A2} predicts the singular set of the
Heun ODE in $\mathbb{Q}(\sqrt{\Delta})$. A natural
strengthening is: \emph{up to PSL$_{2}(\mathbb{Z})$
equivalence, only finitely many distinct CM fields appear
across all degree-two PCFs with bounded coefficients
$|\alpha|,|\beta|,|\gamma|\le M$.} For the sample tested,
six families produce six distinct fields, but
$\mathbb{Q}(\sqrt{-7})$ already appears twice
($\Delta\in\{-7,-28\}$), suggesting the finiteness is
non-trivial.
\label{op:finite-fields}

\item \emph{Borel-resummation channel for the remaining five
families.} The $L(t)$-channel test in
\S\ref{ssec:painleve-deep-channel} cannot recover even
$V_{\mathrm{quad}}$'s known Painlev\'e reduction, indicating
that the Painlev\'e structure for $\Delta<0$ degree-two PCFs
lives in the Borel-resummation Stokes-constant ODE channel
rather than the recurrence-parameter channel used here. A
natural follow-up is to construct a Borel-channel pipeline
and test whether QL01, QL02, QL06, QL15, QL26 admit explicit
Painlev\'e reductions there. The $V_{\mathrm{quad}}$ datum
suggests this is the structurally correct channel, but the
question is open.
\label{op:borel-channel}
\end{enumerate}

%==========================================================
\appendix
\section{Numerical methodology --- detailed}\label{app:numerical}
%==========================================================

\subsection{Wallis recurrence pseudocode}\label{app:wallis}

The reference implementation is given in
Algorithm~\ref{alg:wallis}.

\begin{table}[h]
\centering
\begin{tabular}{p{0.9\textwidth}}
\toprule
\textsc{Algorithm \ref{alg:wallis}: Forward Wallis recurrence with rescaling.}\\
\midrule
\textbf{Input.} Polynomial coefficients
$(\alpha,\beta,\gamma,\delta,\varepsilon,b_{0})$;
target depth $N$; working precision $D$.\\[2pt]
\textbf{Output.} Approximation
$L_{N}=p_{N}/q_{N}$ to $L=K(a_{n},b_{n})$ to
$\ge \rho\cdot N$ digits, where
$\rho=\log_{10}|\Lambda|$.\\[2pt]
1. Set \texttt{mp.dps}~$\leftarrow D$;
$(p_{-1},p_{0},q_{-1},q_{0})\leftarrow(1,b_{0},0,1)$.\\
2. \textbf{For} $n=1,2,\ldots,N$:\\
\quad\quad $a_{n}\leftarrow \delta n+\varepsilon$;
$b_{n}\leftarrow \alpha n^{2}+\beta n+\gamma$.\\
\quad\quad $p_{\mathrm{new}}\leftarrow b_{n}p_{n-1}+a_{n}p_{n-2}$;
$q_{\mathrm{new}}\leftarrow b_{n}q_{n-1}+a_{n}q_{n-2}$.\\
\quad\quad \textbf{If} $|q_{\mathrm{new}}|=0$ \textbf{return} ``diverged''.\\
\quad\quad \textbf{If} $\max(|p_{n-1}|,|q_{n-1}|)>10^{200}$:
rescale $(p_{n-2},p_{n-1},q_{n-2},q_{n-1})$ by maximum
modulus.\\
\quad\quad $(p_{n-2},p_{n-1})\leftarrow(p_{n-1},p_{\mathrm{new}})$;
$(q_{n-2},q_{n-1})\leftarrow(q_{n-1},q_{\mathrm{new}})$.\\
3. \textbf{Return} $p_{N}/q_{N}$.\\
\bottomrule
\end{tabular}
\caption{Forward Wallis recurrence for degree-two PCFs.
Rescaling ensures intermediate operands remain bounded; the
ratio $p_{N}/q_{N}$ is rescale-invariant.}
\label{alg:wallis}
\end{table}

\subsection{Stability digit estimate}\label{app:stab}

After the loop, we estimate the number of stable digits via
\[
\mathrm{stab} \;=\; -\log_{10}|L_{N}-L_{N-5}|,
\]
the ``last six convergents'' gap-closing rate. For all
families in Table~\ref{tab:dichotomy} we observe
$\mathrm{stab}\ge \rho\cdot N$ to within rounding, confirming
the Perron--Poincar\'e exponential decay of the truncation
error.

\subsection{The $F(2,4)$~Trans closed forms}\label{app:closed-forms}

The complete table of $24$ closed forms is reproduced verbatim
from \cite[Appendix~B]{Papanokechi2026F24}; the table is too
long to reproduce here in entirety, but each row is of the
form
\[
L \;=\; \frac{P\big(\pi^{2}, e^{-1}, \log 2, \sqrt{5},
\log\varphi, \arctan(1/\varphi)\big)}{
       Q\big(\pi^{2}, e^{-1}, \log 2, \sqrt{5},
            \log\varphi, \arctan(1/\varphi)\big)}
\]
for $P,Q\in\mathbb{Z}[X_{1},\ldots,X_{6}]$ of degree at most
four in each variable, total degree at most six. PSLQ
detects each row at integer bound $\le 10^{6}$ and working
precision $\le 80$ digits against the eighteen-element basis
$\mathcal{B}_{18}$ of~\eqref{eq:pslq-basis} restricted to its
$F(2,4)$~Trans-relevant sub-basis.

A representative selection of the closed-form rows is given
in Table~\ref{tab:f24-trans-rows}.

\begin{table}[h]
\centering
\small
\begin{tabular}{ccccl}
\toprule
$\alpha$ & $\beta$ & $\gamma$ & $\delta$ & $L = K(a_{n},b_{n})$ \\
\midrule
\multicolumn{5}{l}{\textit{Main class}, $\Delta = +1$, $\Lambda\in\mathbb{Q}$:}\\
$1$ & $\phantom{-}3$ & $1$ & $-2$ & $1/(2/(e-2)-1)$ \\
$1$ & $\phantom{-}3$ & $1$ & $-2$ & $\pi^{2}/6 - 1$ \\
$1$ & $-3$ & $1$ & $-2$ & $1 - 2/\pi^{2}$ \\
$1$ & $\phantom{-}3$ & $1$ & $-2$ & $(2-2/e)/\log 2$ \\
$1$ & $-3$ & $1$ & $-2$ & $1/\zeta(3)$ \\
$1$ & $\phantom{-}3$ & $1$ & $-2$ & $\pi^{2}/(\pi^{2}-6)$ \\
$1$ & $\phantom{-}3$ & $1$ & $-2$ & $(\pi^{2}-6)/(12-\pi^{2})$ \\
\multicolumn{5}{l}{\quad\ldots and $15$ further rows in
the same class, all rational in $\{e^{-1}, \pi^{2}, \log 2,
\zeta(3), \gamma\}$.}\\
\midrule
\multicolumn{5}{l}{\textit{Outlier class},
$\Delta = +8$, $\Lambda = \varphi = (1+\sqrt{5})/2$:}\\
$1$ & $\phantom{-}2$ & $1$ & $+1$ & $\sqrt{5}/(2\log\varphi) - 1/2$ \\
$1$ & $-2$ & $1$ & $+1$ & $\sqrt{5}\,\arctan(1/\varphi) / \log\varphi$ \\
\bottomrule
\end{tabular}
\caption{Representative rows from the $24$-row $F(2,4)$~Trans
closed-form table. Main class ($\Delta=+1$, $22$ rows): all
expressible in $\mathbb{Q}(e^{-1},\pi^{2},\log 2,\zeta(3),\gamma)$.
Outlier class ($\Delta=+8$, $2$ rows): expressible in
$\mathbb{Q}(\sqrt{5})\cdot\{\log\varphi,
\arctan(1/\varphi)\}$. The full table is in
\cite[Appendix~B]{Papanokechi2026F24}; PSLQ detection at
integer bound $\le 10^{6}$ and working precision $\le 80$
digits in every row.}
\label{tab:f24-trans-rows}
\end{table}

\subsection{Painlev\'e ansatz reference forms}\label{app:painleve-forms}

The three Painlev\'e ans\"atze used in the residual matrix
are the standard forms
\begin{align*}
\text{PIII:}\quad
L'' &= \frac{(L')^{2}}{L} - \frac{L'}{t}
+ \frac{AL^{2}+B}{t} + CL^{3} + \frac{D}{L},\\
\text{PV:}\quad
L'' &= \Big(\frac{1}{2L}+\frac{1}{L-1}\Big)(L')^{2}
- \frac{L'}{t}
+ \frac{(L-1)^{2}(AL+B/L)}{t^{2}}+\frac{CL}{t}
+ \frac{DL(L+1)}{L-1},\\
\text{PVI:}\quad
L'' &= \frac{1}{2}\Big(\frac{1}{L}+\frac{1}{L-1}+\frac{1}{L-t}\Big)(L')^{2}
- \Big(\frac{1}{t}+\frac{1}{t-1}+\frac{1}{L-t}\Big)L'\\
&\quad+\frac{L(L-1)(L-t)}{t^{2}(t-1)^{2}}\Big(A
+ \frac{Bt}{L^{2}}
+ \frac{C(t-1)}{(L-1)^{2}}
+ \frac{Dt(t-1)}{(L-t)^{2}}\Big).
\end{align*}
The four parameters $(A,B,C,D)$ are fitted by an
overdetermined linear system at four interior $t$-points and
validated at the two outer points; the maximum validation
residual is reported as the matrix entry. A genuine
identification (HIT) requires residual $\le 10^{-50}$.

\subsection{Per-family residual and Stokes-exponent tables}
\label{app:perfamily}

For each of the four newly probed $\Delta<0$ families
(QL02, QL06, QL15, QL26) we report the full
$2\times 3$ Painlev\'e residual matrix and the
$2\times 6$ Stokes-exponent matrix in
Tables~\ref{tab:full-residuals}--\ref{tab:full-S}.
Cells in italics fall below the AMBIG threshold
$10^{-20}$; no cell does. The $5$-pt-vs-$7$-pt
stencil disagreement (column ``stencil'') is reported as
the maximum across both deformations and all four
derivatives.

\begin{table}[h]
\centering
\small
\begin{tabular}{lllrrrr}
\toprule
Family & $\Delta$ & Def. & P-III  & P-V  & P-VI  & stencil \\
\midrule
QL02 & $-4$  & A & $9.39\times 10^{-4}$  & $9.64\times 10^{-3}$ & $0.129$    & $2.36\times 10^{-7}$ \\
QL02 & $-4$  & B & $6.07\times 10^{-2}$  & $4.83\times 10^{-2}$ & $1.62$     & $2.60\times 10^{-5}$ \\
QL06 & $-7$  & A & $1.25\times 10^{-4}$  & $1.07\times 10^{-4}$ & $2.05\times 10^{-2}$  & $\le 10^{-7}$ \\
QL06 & $-7$  & B & $8.82\times 10^{-3}$  & $0.133$    & $5.45$     & $\le 10^{-5}$ \\
QL15 & $-20$ & A & $1.87\times 10^{-2}$  & $0.230$    & $0.450$    & $\le 10^{-7}$ \\
QL15 & $-20$ & B & $0.214$    & $0.301$    & $7.10$     & $\le 10^{-5}$ \\
QL26 & $-28$ & A & $2.40\times 10^{-2}$  & $2.49\times 10^{-2}$ & $36.3$     & $\le 10^{-7}$ \\
QL26 & $-28$ & B & $4.79\times 10^{-2}$  & $0.884$    & $315$      & $\le 10^{-5}$ \\
\bottomrule
\end{tabular}
\caption{Painlev\'e residual matrices for the four newly
probed $\Delta<0$ families. No cell achieves the HIT
threshold $10^{-50}$ or the AMBIG threshold $10^{-20}$.
Residuals reported are the maximum validation residuals at
the two outer points $t=\pm 0.3$ after fitting at
$t\in\{\pm 0.1,\pm 0.2\}$. Working precision
$\mathrm{dps}=250$, depth $1500$.}
\label{tab:full-residuals}
\end{table}

\begin{table}[h]
\centering
\small
\begin{tabular}{llrrrrrr}
\toprule
Family & Def. & $S(+0.1)$ & $S(-0.1)$ & $S(+0.2)$ & $S(-0.2)$ & $S(+0.3)$ & $S(-0.3)$ \\
\midrule
QL02 & A & $1.155$  & $1.165$  & $1.215$  & $1.245$  & $1.279$  & $1.339$ \\
QL02 & B & $0.828$  & $0.892$  & $0.710$  & $0.894$  & $\mathbf{0.558}$ & $0.926$ \\
QL06 & A & $1.025$  & $1.026$  & $1.034$  & $1.038$  & $1.045$  & $1.052$ \\
QL06 & B & $0.713$  & $0.763$  & $0.556$  & $0.700$  & $\mathbf{0.364}$ & $0.655$ \\
QL15 & A & $0.943$  & $0.949$  & $0.838$  & $0.871$  & $0.708$  & $0.785$ \\
QL15 & B & $1.144$  & $1.196$  & $1.171$  & $1.299$  & $1.187$  & $1.398$ \\
QL26 & A & $0.639$  & $0.633$  & $0.488$  & $0.469$  & $0.321$  & $0.282$ \\
QL26 & B & $0.345$  & $0.389$  & $0.036$  & $0.160$  & $\mathbf{-0.324}$ & $-0.072$ \\
\bottomrule
\end{tabular}
\caption{Stokes-exponent $S(t)$ values
for the four newly probed $\Delta<0$ families. The minimum
across the table per family (in bold) gives $S_{\min}$ in
Table~\ref{tab:6family}. $S<1$ on at least one deformation
arc in every family; QL26 reaches a negative
$S(+0.3)=-0.324$ on the D-B arc, the strongest signal in
the database. Note that $S(t)>1$ entries are also
diagnostically informative: a single deformation that
\emph{lifts} $S$ above $1$ does not refute
Conjecture~\ref{conj:A3}, which only requires \emph{some}
CM-respecting deformation to give $S<1$.}
\label{tab:full-S}
\end{table}

The QL15 family is unusual in that the D-A arc
(constant-term scaling) gives $S<1$ uniformly, while the
D-B arc (root-radius scaling) gives $S>1$ uniformly.
This is the only family in the sample with this asymmetry,
and it indicates that the deformation choice is not
neutral: the structural Stokes signal lives on the
$\gamma$-shift arc for QL15, not on the radius-scaling arc.
The reverse is true for QL02, QL06, and QL26, where the
D-B arc provides the strongest signal.

\subsection{Reproducibility URLs}\label{app:repro}

The session-level computation scripts and per-family
residual JSON files are archived at the SIARC research
workspace under
\begin{itemize}
\item \texttt{sessions/2026-04-30/QL15-STOKES/} (QL15 baseline);
\item \texttt{sessions/2026-05-01/STOKES-FAMILY-EXTEND-3X/}
(QL02, QL06, QL26 extensions);
\end{itemize}
together with the shared engine
\texttt{stokes\_session\_b.py} that runs the six-step
Session~B template given a family-spec dictionary. The
$F(2,4)$~Trans closed-form database is at
\texttt{sessions/2026-04-13/F24-TRANS-TABLE/}. Files are
available from the author on request, in line with the AEAL
reproducibility requirement.
residual JSON files are archived at the SIARC research
workspace under
\begin{itemize}
\item \texttt{sessions/2026-04-30/QL15-STOKES/} (QL15 baseline);
\item \texttt{sessions/2026-05-01/STOKES-FAMILY-EXTEND-3X/}
(QL02, QL06, QL26 extensions);
\end{itemize}
together with the shared engine
\texttt{stokes\_session\_b.py} that runs the six-step
Session~B template given a family-spec dictionary. The
$F(2,4)$~Trans closed-form database is at
\texttt{sessions/2026-04-13/F24-TRANS-TABLE/}. Files are
available from the author on request, in line with the AEAL
reproducibility requirement.

%==========================================================
\begin{thebibliography}{99}

\bibitem{Papanokechi2026Spec}
Papanokechi, \emph{Spectral classes of polynomial continued
fractions: a hypergeometric framework and arithmetic
obstructions}.
Preprint, 2026.

\bibitem{Papanokechi2026Vquad}
Papanokechi, \emph{Resurgent analysis of the
$V_{\mathrm{quad}}$ continued fraction}.
Zenodo preprints, 2026.
\href{https://doi.org/10.5281/zenodo.19491767}
{DOI:10.5281/zenodo.19491767},
\href{https://doi.org/10.5281/zenodo.19491768}
{DOI:10.5281/zenodo.19491768}.

\bibitem{Papanokechi2026F24}
Papanokechi, \emph{Arithmetic stratification of the
$F(2,4)$ family of polynomial continued fractions}.
Preprint, 2026.

\bibitem{P7AIBehSci}
Papanokechi, \emph{AI behavior science: a proposal for an
empirical discipline} (P7). Zenodo, 2026.
\href{https://doi.org/10.5281/zenodo.19562751}
{DOI:10.5281/zenodo.19562751}.

\bibitem{P8AEAL}
Papanokechi, \emph{AEAL: AI-augmented evidence aggregation
and logging --- a multi-agent methodology for empirical
mathematics} (P8). Zenodo, 2026.
\href{https://doi.org/10.5281/zenodo.19565086}
{DOI:10.5281/zenodo.19565086}.

\bibitem{Okamoto1987}
K.~Okamoto, \emph{Studies on the Painlev\'e equations IV: third
Painlev\'e equation $P_{III}$}. Funkcial.\ Ekvac.\ \textbf{30}
(1987), 305--332.

\bibitem{FergusonBailey1992}
H.~R.~P.~Ferguson and D.~H.~Bailey, \emph{A polynomial-time,
numerically stable integer relation algorithm}. RNR Tech.\ Rep.\
RNR-91-032, NASA Ames, 1992.

\bibitem{Raayoni2021}
G.~Raayoni, S.~Gottlieb, Y.~Manor, G.~Pisha, Y.~Harris,
U.~Mendlovic, D.~Haviv, Y.~Hadad, and I.~Kaminer,
\emph{Generating conjectures on fundamental constants with
the Ramanujan Machine}. Nature \textbf{590} (2021), 67--73.

\bibitem{Schneider1934}
T.~Schneider, \emph{Transzendenzuntersuchungen periodischer
Funktionen}. J.\ reine angew.\ Math.\ \textbf{172} (1934),
65--69.

\bibitem{Chudnovsky1980}
G.~V.~Chudnovsky, \emph{Algebraic independence of constants
connected with exponential and elliptic functions}.
Dokl.\ Akad.\ Nauk Ukrain.\ SSR (1980), 8--12.

\bibitem{Wolfart1988}
J.~Wolfart, \emph{Werte hypergeometrischer Funktionen}.
Invent.\ Math.\ \textbf{92} (1988), 187--216.

\bibitem{Bertrand2002}
D.~Bertrand, \emph{Theta functions and transcendence}.
Ramanujan~J.\ \textbf{1} (2002), 339--350.

\bibitem{HeunRef}
A.~Ronveaux (ed.), \emph{Heun's Differential Equations}.
Oxford Univ.\ Press, 1995.

\bibitem{IwasakiKimuraShimomuraYoshida}
K.~Iwasaki, H.~Kimura, S.~Shimomura, M.~Yoshida,
\emph{From Gauss to Painlev\'e: a Modern Theory of Special
Functions}. Vieweg, 1991.

\bibitem{Jimbo1982}
M.~Jimbo, \emph{Monodromy problem and the boundary
condition for some Painlev\'e equations}.
Publ.\ Res.\ Inst.\ Math.\ Sci.\ \textbf{18} (1982), 1137--1161.

\bibitem{Boutroux1913}
P.~Boutroux, \emph{Recherches sur les transcendantes de M.\
Painlev\'e}. Ann.\ Sci.\ \'Ec.\ Norm.\ Sup\'er.\ \textbf{30}
(1913), 255--375.

\bibitem{EcalleResurgence}
J.~\'Ecalle, \emph{Les fonctions r\'esurgentes}, vols.\ I--III.
Publ.\ Math.\ d'Orsay, 1981--1985.

\bibitem{LefschetzBessis}
M.~V.~Berry and C.~J.~Howls, \emph{Hyperasymptotics for
integrals with saddles}.
Proc.\ R.\ Soc.\ Lond.\ A \textbf{434} (1991), 657--675.

\bibitem{mpmath}
F.~Johansson et al.,
\emph{mpmath: a Python library for arbitrary-precision
floating-point arithmetic} (version 1.3.0), 2023.
\url{http://mpmath.org}.

\bibitem{Stieltjes1894}
T.~J.~Stieltjes, \emph{Recherches sur les fractions continues}.
Ann.\ Fac.\ Sci.\ Toulouse \textbf{8} (1894), J1--J122.

\bibitem{LangElliptic}
S.~Lang, \emph{Elliptic Functions}, 2nd ed.\ Graduate Texts
in Mathematics \textbf{112}, Springer, 1987.

\bibitem{SilvermanAdvanced}
J.~H.~Silverman, \emph{Advanced Topics in the Arithmetic of
Elliptic Curves}. Graduate Texts in Mathematics \textbf{151},
Springer, 1994.

\bibitem{pcfdatabase2026}
Papanokechi, \emph{PCF spectral research database
(first batch)}. Computational artifact, 2026, available
on request.

\end{thebibliography}

\end{document}
